\SourceLeaf{183}of the complete inconclusiveness of the matter, should that
proof against Euclidean geometry turn out, contrary to what I could have expected,
to lack rigor --- [had been recognized by me], and partly because
the many earlier investigations in plane geometry would still be
worth keeping and could be used with a few modifications,
and because, at least up to a fairly wide limit that might perhaps
be determined more precisely, the propositions of solid
geometry and mechanics would be approximately valid, I found this
evening --- while trying to find a way into your
transcendental trigonometry and turning, because I was unable
to obtain adequate, determinate functions for it in the plane,
to constructions in space, much to my
surprise, the following \textls{proof for the Euclidean theory of parallels}. .\ .\ .\ .\ .

.\ .\ .\ .\ .\ .\ Just as I am about to close, I notice that the foregoing
proof for the Euclidean theory of parallels is faulty. .\ .\ .\ .\ . The hope of
reaching a fully decided result would thus have disappeared here too,
and I must again rest content with what I stated earlier.
I do believe, however, that along that path I have taken at least one step toward your
transcendental trigonometry, since with the aid of
spherical trigonometry I can give the ratios of all the constants, at least
\textls{by constructing right-angled} triangles. I still lack
the actual calculation of the base of an isosceles triangle from
its side. For this I shall try to start from the equilateral triangle.
.\ .\ .\}

\begin{center}\rule{3cm}{0.4pt}\end{center}

\Needspace{11\baselineskip}
\begin{center}\textsc{Remark.}\end{center}

In the letter of 12 December 1816, extracts of which are given here,
\textsc{Wachter} repeatedly refers to a conversation with \textsc{Gauss},
which took place during his »last stay in Göttingen« and concerned
anti-Euclidean geometry. A letter from
\textsc{Gauss} to \textsc{Olbers} of 4 June 1816 shows that this
visit by \textsc{Wachter} took place in April 1816.

\begin{flushright}\textsc{Stäckel.}\end{flushright}

\begin{center}\rule{3cm}{0.4pt}\end{center}

